\def\thoigian{2}
\section
[Giải bài toán bằng cách lập hệ phương trình]
{GIẢI BÀI TOÁN BẰNG CÁCH LẬP HỆ PHƯƠNG TRÌNH}

\subsection{MỤC TIÊU}

\paragraph{Về kiến thức, kĩ năng}

\begin{itemize}
	\item Giải một số bài toán bằng cách lập hệ hai phương trình bậc nhất hai ẩn.
\end{itemize}

\paragraph{Về năng lực}

\begin{itemize}
	\item Rèn luyện năng lực toán học, nói riêng là năng lực mô hình hoá toán học và năng lực giải quyết vấn đề toán học.
	\item Bồi dưỡng hứng thú học tập, ý thức làm việc nhóm, ý thức tìm tòi, khám phá và sáng tạo cho HS. 
\end{itemize}

\paragraph{Về phẩm chất}

\begin{itemize}
	\item Góp phần giúp HS rèn luyện và phát triển các phẩm chất tốt đẹp (yêu nước, nhân ái, chăm chỉ, trung thực, trách nhiệm):
	\begin{itemize}
		\item Tích cực phát biểu, xây dựng bài và tham gia các hoạt động nhóm;
		\item Có ý thức tích cực tìm tòi, sáng tạo trong học tập; phát huy điểm mạnh, khắc phục các điểm yếu của bản thân. 
	\end{itemize}
\end{itemize}

\subsection{THIẾT BỊ DẠY HỌC VÀ HỌC LIỆU}

\paragraph{Giáo viên}

\begin{itemize}
	\item Giáo án, bảng phụ, máy chiếu (nếu có),… 
\end{itemize}

\paragraph{Học sinh}

\begin{itemize}
	\item SGK, vở ghi, dụng cụ học tập.
\end{itemize}

\subsection{TIẾN TRÌNH DẠY HỌC}

\danhsachtiet

\subsubsection{Giải bài toán bằng cách lập hệ phương trình}

\paragraph{Hoạt động khởi động}

\muctieu

\begin{itemize}
	\item Gợi động cơ, tạo tình huống xuất hiện trong thực tế để HS vận dụng kiến thức về giải hệ phương trình để giải quyết tính huống.
\end{itemize}

\noidungsanpham

\begin{modau}
	Một vật có khối lượng $124$ g và thể tích $15$ cm$^3$ là hợp kim của đồng và kẽm. TÍnh xem trong đó có bao nhiêu gam đồng và bao nhiêu gam kẽm, biết rằng $1$ cm$^3$ đồng nặng $8{,}9$ g và $1$ cm$^3$ kẽm nặng $7$ g.
\end{modau}

\tochucthuchien

\begin{tcth}
	\buocmot
	\item GV yêu cầu HS đọc tình huống thực tế và cho biết có bao nhiêu đại lượng chưa biết trong bài \textit{(GV giải thích lại thế nào là đại lượng nếu cần)}.
	\buochai
	\item HS suy nghĩ về tình huống mở đầu và nảy sinh nhu cầu tìm hiểu cách giải quyết tình hướng này.
	\buocba
	\item HS trả lời câu hỏi (không yêu cầu trả lời chính xác).
	\buocbon
	\item GV giới thiệu bài học mới về giải bài toán bằn cách lập hệ phương trình.
\end{tcth}

\paragraph{Hoạt động hình thành kiến thức}

\muctieu

\begin{itemize}
	\item HS nhận biết các bước thực hiện khi giải bài toán bằng cách lập hệ phương trình.
\end{itemize}

\noidungsanpham

\begin{ttkp}
	\textbf{Giải bài toàn bằng cách lập hệ phương trình}\\
	Xét bài toán ở Tình huống mở đầu. Gọi $x$ là số gam đồng, $y$ là số gam kẽm cần tính.
	\begin{itemize}
		\item[] \textbf{HĐ1:} Biểu thị khối lượng của vật qua $x$ và $y$.
		\item[] \textbf{HĐ2:} Biểu thị thể tích của vật qua $x$ và $y$.
		\item[] \textbf{HĐ3:} Giải hệ gồm hai phương trình bậc nhất hai ẩn $x$, $y$ nhận được ở \textbf{HĐ1} và \textbf{HĐ2}. Từ đó trả lời câu hỏi ở Tình huống mở đầu.
	\end{itemize}
	%%%
	\doicot
	%%%
	\huongdan
	\begin{itemize}
		\item[] \textbf{HĐ1:} Khối lượng của vật là $x+y$ (g).\\
		Mà khối lượng của vật đồng thời là $124$ g, nên 
		$$x+y=124\qquad(1)$$
		\item[] \textbf{HĐ2:} Thể tích của vật qua $\dfrac{x}{8{,}9}+\dfrac{y}{7}$ (cm$^3$).\\
		Mà thể tích của vật đồng thời là $15$ cm$^3$, nên 
		$$\dfrac{x}{8{,}9}+\dfrac{y}{7}=15\qquad(2)$$
		\item[] \textbf{HĐ3:} Từ (1) và (2), ta lập được hệ phương trình
		$$\heva{&x+y=124\\&\dfrac{x}{8{,}9}+\dfrac{y}{7}=15.}$$
		Giải hệ phương trình trên ta được:
		$$x=89;y=35$$
	\end{itemize}
\end{ttkp}

\begin{vidu}
	Tìm hai số tự nhiên có tổng bằng $1006$, biết rằng nếu lấy số lớn chia cho số nhỏ thì được thương là $2$ và số dư là $124$.
	%%%
	\doicot
	%%%
	\begin{itemize}
		\item Gọi hai số cần tìm là $x$ và $y$, trong đó $x < y$. Số dư trong phép chia $y$ cho $x$ là $124$ nên $x > 124$. Vậy điều kiện của hai ẩn là $x, y \in \mathbb{N}$ và $124 < x < y$.\\
		Tổng hai số bằng $1\,006$ nên ta có phương trình $x + y = 1\,006$.\\
		Khi chia $y$ cho $x$ ta được thương là $2$, dư $124$ nên ta có phương trình $y = 2x + 124$.\\
		Do đó, ta có hệ phương trình 
		$$\heva{&x + y = 1\,006 & (1) \\& y = 2x + 124. & (2) }$$
		\item Giải hệ phương trình: \\
		Từ (2) thế $y = 2x + 124$ vào (1), ta được $3x + 124 = 1\,006$ hay $3x = 882$, suy ra $x = 294$. \\
		Từ đó ta được $y = 2 \cdot 294 + 124 = 712$.
		\item Các giá trị $x = 294$ và $y = 712$ thoả mãn các điều kiện của ẩn. \\
		Vậy hai số cần tìm là $294$ và $712$.
	\end{itemize}
\end{vidu}

\tochucthuchien

\begin{tcth}[ttkp]
	\buocmot
	\item GV cho HS hoạt động cá nhân trong 5 phút. Sau đó, GV gọi một HS lên bảng trả lời hai HĐ1, 2; các HS khác quan sát, nhận xét và góp ý phần lời giải của HS. Tiếp theo GV gọi một HS khác lên làm HĐ3. 
	\buochai
	\item HS thực hiện theo hướng dẫn của GV.
	\buocba
	\item GV tổng kết suy ra cách giải bài toán bằng cách lập hệ phương trình.
	\buocbon
	\item GV viết bảng hoặc trình chiếu Nhận xét:
	\begin{nhanxet}
		Các bước giải một bài toán bằng cách lập hệ phương trình:
		\begin{cacbuoc}
			\item Lập hệ phương trình:
			\begin{itemize}
				\item Chọn ẩn số (thường chọn hai ẩn số) và đặt điều kiện thích hợp cho các ẩn số;
				\item Biểu diễn các đại lượng chưa biết theo ẩn và các đại lượng đã biết;
				\item Lập hệ phương trình biểu thị mối quan hệ giữa các đại lượng.
			\end{itemize}
			\item Giải hệ phương trình.
			\item Trả lời: Kiểm tra xem trong các nghiệm tìm được của hệ phương trình, nghiệm nào thoả mãn, nghiệm nào không thoả mãn điều kiện của ẩn, rồi kết luận.
		\end{cacbuoc}
	\end{nhanxet}
\end{tcth}

\begin{tcth}[vidu=1]
	\buocmot
	\item GV hướng dẫn HS lần lượt thực hiện theo ba Bước 1, 2, 3 nêu trong Nhận xét ở trên.
	\buochai
	\item HS làm bài vào vở dưới sự hướng dẫn của GV.
	\item GV gọi HS lên bảng làm bài.
	\buocba
	\item Các HS khác nhận xét. 
	\buocbon
	\item GV nhận xét, chỉnh sửa nếu cần.
\end{tcth}

\paragraph{Hoạt động luyện tập}

\muctieu

\begin{itemize}
	\item Củng cố kĩ năng giải bài toán bằng cách lập hệ phương trình.
\end{itemize}

\noidungsanpham

\begin{luyentap}
	Một chiếc xe khách đi từ Thành phố Hồ Chí Minh đến Cần Thơ, quãng đường dài $170\text{ km}$. Sau khi xe khách xuất phát $1$ giờ $40$ phút, một chiếc xe tải bắt đầu đi từ Cần Thơ về Thành phố Hồ Chí Minh (trên cùng một tuyến đường với xe khách) và gặp xe khách sau đó $40$ phút. Tính vận tốc của mỗi xe, biết rằng mỗi giờ xe khách đi nhanh hơn xe tải $15\text{ km}$.
	%%%
	\doicot
	%%%
	\huongdan\\
	Gọi $x\text{ (km/h)}$ là vận tốc của xe tải và $y\text{ (km/h)}$ là vận tốc xe khách ($x, y > 0$). Chú ý rằng hai xe (đi ngược chiều) gặp nhau khi tổng quãng đường hai xe đã đi bằng $170\text{ km}$.
\end{luyentap}

\begin{vidu}
	Hai đội công nhân cùng làm một đoạn đường trong $24$ ngày thì xong. Mỗi ngày, đội I làm được nhiều gấp rưỡi đội II. Hỏi nếu làm một mình thì mỗi đội làm xong đoạn đường đó trong bao lâu? (Giả sử năng suất của mỗi đội là không đổi).
	%%%
	\doicot
	%%%
	\begin{itemize}
		\item Gọi $x$ là số ngày để đội I hoàn thành công việc nếu làm một mình; $y$ là số ngày để đội II hoàn thành công việc nếu làm một mình. Điều kiện: $x > 0$ và $y > 0$.\\
		Mỗi ngày đội I làm được $\dfrac{1}{x}$ (công việc) và đội II làm được $\dfrac{1}{y}$ (công việc).\\
		Mỗi ngày đội I làm được nhiều gấp rưỡi đội II nên ta có phương trình $\dfrac{1}{x} = 1,5 \cdot \dfrac{1}{y}$ hay 
		$$\dfrac{1}{x} = \dfrac{3}{2} \cdot \dfrac{1}{y} \quad (1)$$
		Hai đội làm chung trong $24$ ngày thì xong công việc nên mỗi ngày, hai đội làm chung thì được $\dfrac{1}{24}$ (công việc). Ta có phương trình
		$$\dfrac{1}{x} + \dfrac{1}{y} = \dfrac{1}{24} \quad (2)$$
		Từ (1) và (2), ta có hệ phương trình 
		$$(I)\quad\heva{&\dfrac{1}{x} = \dfrac{3}{2} \cdot \dfrac{1}{y} \\ &\dfrac{1}{x} + \dfrac{1}{y} = \dfrac{1}{24}.}$$
		\item Nếu đặt $u = \dfrac{1}{x}$ và $v = \dfrac{1}{y}$ thì ta có hệ phương trình bậc nhất hai ẩn mới là $u$ và $v$:
		$$(II)\quad\heva{&u = \dfrac{3}{2}v & (3) \\ &u + v = \dfrac{1}{24}. & (4)}$$
		Giải hệ (II): Thế $u = \dfrac{3}{2}v$ vào phương trình (4), ta được $\dfrac{3}{2}v + v = \dfrac{1}{24}$, hay $\dfrac{5}{2}v = \dfrac{1}{24}$, suy ra $v = \dfrac{1}{60}$. Do đó $u = \dfrac{3}{2}v = \dfrac{3}{2} \cdot \dfrac{1}{60} = \dfrac{1}{40}$.\\
		Từ đó, ta có:
		$u = \dfrac{1}{x} = \dfrac{1}{40}$ suy ra $x = 40$; $v = \dfrac{1}{y} = \dfrac{1}{60}$ suy ra $y = 60$.
		\item Các giá trị tìm được của $x$ và $y$ thoả mãn điều kiện của ẩn.
	\end{itemize}
	\textit{Trả lời:} Nếu làm một mình thì đội I làm xong đoạn đường đó trong $40$ ngày, còn đội II làm xong trong $60$ ngày.
\end{vidu}

\begin{luyentap}
	Nếu hai vòi nước cùng chảy vào một bể không có nước thì bể sẽ đầy trong 1 giờ 20 phút. Nếu mở riêng vòi thứ nhất trong 10 phút và vòi thứ hai trong 12 phút thì chỉ được $\dfrac{2}{15}$ bể nước. Hỏi nếu mở riêng từng vòi thì thời gian để mỗi vòi chảy đầy bể nước là bao nhiêu phút?
	%%%
	\doicot
	%%%
	\huongdan\\
	Gọi $x$ (phút là thời gian vòi thứ nhất chảy đầy bể; $y$ (phút) là thời gian vòi thứ hai chảy đầy bể.\\
	Hệ phương trình:
	$$\heva{&80\left(\dfrac{1}{x}+\dfrac{1}{y}\right)=1\\&\dfrac{10}{x}+\dfrac{12}{y}=\dfrac{2}{15}.}$$
	Vòi thứ nhất: 120 phút, vòi thứ hai: 240 phút.
\end{luyentap}

\tochucthuchien

\begin{tcth}[luyentap=1]
	\buocmot
	\item GV cho HS hoạt động cá nhân.
	\buochai
	\item HS làm việc dưới sự hướng dẫn của GV.
	\item GV gọi HS lên bảng làm bài.
	\buocba
	\item Các HS khác nhận xét. 
	\buocbon
	\item GV nhận xét, chỉnh sửa nếu cần.
\end{tcth}

\begin{tcth}[vidu=2]
	\buocmot
	\item GV cần giải thích tương quan sau cho HS: ``Nếu đơn vị $A$ làm xong công việc (coi là $1$ công việc) trong $n$ ngày, thì mỗi ngày đơn vị $A$ làm $\dfrac{1}{n}$ công việc''.
	\item GV cho HS hoạt động cá nhân.
	\buochai
	\item HS làm việc dưới sự hướng dẫn của GV.
	\item GV gọi HS lên bảng làm bài.
	\buocba
	\item Các HS khác nhận xét. 
	\buocbon
	\item GV nhận xét, chỉnh sửa nếu cần.
	\item Lưu ý cho GV:\\
	\textit{Phương pháp đặt ẩn phụ không được trình bày tường minh trong SGK về mặt lí thuyết, GV cần lưu ý cho HS là tuy hệ (I) không phải là hệ bậc nhất hai ẩn, nhưng nếu ta đặt (ẩn phụ) $u=\dfrac{1}{x}$, $v=\dfrac{1}{y}$ thì ta lại được một hệ bậc nhất hai ẩn (II) đối với hai ẩn mới là $u$, $v$. Hệ này đơn giản hơn hệ (I) và HS đã biết cách giải.}
\end{tcth}

\begin{tcth}[luyentap=2]
	\buocmot
	\item GV cho HS hoạt động theo nhóm đôi.
	\buochai
	\item HS trao đổi nhóm để hoàn thành yêu cầu của Luyện tập 2.
	\item GV gọi 1 HS lên bảng làm bài.
	\buocba
	\item Các HS khác nhận xét. 
	\buocbon
	\item GV nhận xét, chỉnh sửa nếu cần.
\end{tcth}

\paragraph{Tổng kết và hướng dẫn công việc ở nhà}

GV tổng kết lại nội dung bài học và dặn dò công việc ở nhà cho HS:

\begin{itemize}
	\item GV tổng kết lại các kiến thức trọng tâm của bài học: Giải bài toán bằng cách lập hệ phương trình.
	\item Nhắc HS về nhà ôn tập các nội dung đã học.
	\item Giao cho HS làm các bài tập sau trong SGK: Bài 1.15 đến Bài 1.18.
\end{itemize}

\subsubsection{Chữa bài tập}

\paragraph{Hoạt động khởi động}

\muctieu

\begin{itemize}
	\item HS nhớ lại các cách giải bài toán bằng cách lập hệ phương trình bậc nhất hai ẩn.
\end{itemize}

\noidungsanpham

\begin{ontap}
	Nhắc lại các bước thực hiện khi giải bài toán bằng cách lập hệ phương trình.
	%%%
	\doicot
	%%%
	Các bước giải một bài toán bằng cách lập hệ phương trình:
	\begin{cacbuoc}
		\item Lập hệ phương trình:
		\begin{itemize}
			\item Chọn ẩn số (thường chọn hai ẩn số) và đặt điều kiện thích hợp cho các ẩn số;
			\item Biểu diễn các đại lượng chưa biết theo ẩn và các đại lượng đã biết;
			\item Lập hệ phương trình biểu thị mối quan hệ giữa các đại lượng.
		\end{itemize}
		\item Giải hệ phương trình.
		\item Trả lời: Kiểm tra xem trong các nghiệm tìm được của hệ phương trình, nghiệm nào thoả mãn, nghiệm nào không thoả mãn điều kiện của ẩn, rồi kết luận.
	\end{cacbuoc}
\end{ontap}

\tochucthuchien

\begin{tcth}
	\buocmot
	\item GV yêu cầu nhắc lại các bước thực	hiện khi giải bài toán bằng cách lập hệ	phương trình.
	\buochai
	\item HS chuẩn bị câu trả lời
	\buocba
	\item HS giơ tay phát biểu.
	\item Các HS khác nhận xét.
	\buocbon
	\item GV trình chiếu lại nội dung kiến thức và nhắc lại các bước cho HS.
\end{tcth}

\paragraph{Hoạt động luyện tập}

\muctieu

\begin{itemize}
	\item Củng cố kĩ năng giải bài toán bằng cách lập hệ phương trình.
\end{itemize}

\noidungsanpham

\begin{baitap}[1.15]
	Tìm số tự nhiên $n$ có hai chữ số, biết rằng tổng của hai chữ số đó bằng $12$, và nếu viết hai chữ số đó theo thứ tự ngược lại thì được một số lớn hơn $n$ là $36$ đơn vị.
	%%%
	\doicot
	%%%
	\huongdan\\
	Gọi $x$ là chữ số hàng chục, $y$ là chữ số hàng trăm (khi đó $n = 10x + y$). Điều kiện của ẩn là: $x, y \in \mathbb{N}$ và $0 < x \le 9$ và $0 \le y \le 9$. Ta có phương trình thứ nhất là $x + y = 12$. Khi viết hai chữ số của $n$ theo thứ tự ngược lại, ta được số $10y + x$.\\	
	Theo giả thiết ta có phương trình 
	$$(10y + x) - (10x + y) = 36,$$ 
	hay $9y - 9x = 36$.\\
	Từ đó ta có hệ phương trình
	$$\heva{&x+y=12\\&9y-9x=36.}$$
\end{baitap}

\begin{baitap}[1.16]
	Điểm số trung bình của một vận động viên bắn súng sau $100$ lần bắn là $8,69$ điểm. Kết quả cụ thể được ghi trong bảng sau, trong đó có hai ô bị mờ không đọc được (đánh dấu "?"):
	\begin{center}
		\begin{tblr}{|C{3cm}|c|c|c|c|c|}
			\hline
			Điểm số của mỗi lần bắn & 10 & 9 & 8 & 7 & 6 \\ \hline
			Số lần bắn & 25 & 42 & ? & 15 & ? \\ \hline
		\end{tblr}
	\end{center}
	Em hãy tìm lại các số bị mờ trong hai ô đó.
	%%%
	\doicot
	%%%
	\huongdan\\
	Gọi $x$ là số lần bắn được 8 điểm và $y$ là số lần bắn được 6 điểm. Điều kiện là $x, y \in \mathbb{N}$ (cũng có thể nêu điều kiện $x, y < 100$). Theo giả thiết, ta có hệ phương trình
	$$\heva{&25 + 42 + x + 15 + y = 100 \\ &(10 \cdot 25 + 9 \cdot 42 + 8x + 7 \cdot 15 + 6y) : 100 = 8{,}69}$$
	hay
	$$\heva{&x + y = 18 \\ &8x + 6y = 136.}$$
\end{baitap}

\begin{baitap}[1.17]
	Năm ngoái, hai đơn vị sản xuất nông nghiệp thu hoạch được $3\,600$ tấn thóc. Năm nay, hai đơn vị thu hoạch được $4\,095$ tấn thóc. Hỏi năm nay, mỗi đơn vị thu hoạch được bao nhiêu tấn thóc, biết rằng năm nay, đơn vị thứ nhất làm vượt mức $15\%$, đơn vị thứ hai làm vượt mức $12\%$ so với năm ngoái?\\	
	Hãy dùng máy tính cầm tay để kiểm tra lại kết quả thu được.
	%%%
	\doicot
	%%%
	\huongdan\\
	Gọi $x$ và $y$ lần lượt là số tấn thóc mà đơn vị thứ nhất và thứ hai thu hoạch được trong năm ngoái ($x, y > 0$). Ta có hệ phương trình
	$$\heva{&x + y = 3600 \\ &115\%x + 112\%y = 4095}$$
	hay
	$$\heva{&x + y = 3600 \\ &1,15x + 1,12y = 4095.}$$
\end{baitap}

\begin{baitap}[1.18]
	Hai người thợ cùng làm một công việc trong $16$ giờ thì xong. Nếu người thứ nhất làm trong $3$ giờ và người thứ hai làm trong $6$ giờ thì chỉ hoàn thành được $25\%$ công việc. Hỏi nếu làm riêng thì mỗi người hoàn thành công việc trong bao lâu?
	%%%
	\doicot
	%%%
	Gọi $x$ (giờ) là thời gian để người thứ nhất hoàn thành công việc, $y$ (giờ) là thời gian để người thứ hai hoàn thành công việc (điều kiện: $x > 0, y > 0$).\\
	Khi đó người thứ nhất mỗi giờ hoàn thành được $\dfrac{1}{x}$ công việc; người thứ hai được $\dfrac{1}{y}$ công việc.\\
	Cả hai cùng làm thì mỗi giờ được $\dfrac{1}{x} + \dfrac{1}{y}$ công việc, và hoàn thành toàn bộ công việc trong $16$ giờ nên ta có phương trình $16 \left( \dfrac{1}{x} + \dfrac{1}{y} \right) = 1$.\\
	Người thứ nhất làm trong $3$ giờ được $\dfrac{3}{x}$ công việc; người thứ hai làm trong $6$ giờ được $\dfrac{6}{y}$ công việc và khi đó cả hai chỉ hoàn thành được $25\%$ ($= \dfrac{1}{4}$) công việc nên ta có phương trình $\dfrac{3}{x} + \dfrac{6}{y} = \dfrac{1}{4}$.\\
	Ta có hệ phương trình (I) $\heva{&16 \left( \dfrac{1}{x} + \dfrac{1}{y} \right) = 1 \\ &\dfrac{3}{x} + \dfrac{6}{y} = \dfrac{1}{4}.}$\\
	Bằng cách đặt $u = \dfrac{1}{x}$ và $v = \dfrac{1}{y}$, ta đưa $(I)$ về dạng: (II) $\heva{&16(u + v) = 1 & (1) \\ &3u + 6v = \dfrac{1}{4} & (2)}$\\
	Giải hệ phương trình trên, ta được: Nếu làm riêng thì người thứ nhất hoàn thành công việc trong $24$ giờ, người thứ hai trong $48$ giờ.
\end{baitap}

\tochucthuchien

\begin{tcth}[baitap={từ 1.15 đến 1.18}]
	\buocmot
	\item GV cho HS làm bài cá nhân.
	\buochai
	\item HS làm bài vào vở.
	\item GV gọi 1 HS lên bảng làm bài.
	\buocba
	\item Các HS khác nhận xét. 
	\buocbon
	\item GV nhận xét, chỉnh sửa nếu cần.
\end{tcth}







